\documentclass{article}
 
\usepackage{booktabs}
\usepackage{tabularx}
\usepackage{array}
\usepackage{caption}
\usepackage[a4paper, margin=2cm]{geometry}
 
\begin{document}
 
\begin{table}[ht]
\centering
 
\caption*{\textbf{Dicionário de dados}}
 
\vspace{0.2cm}
 
\renewcommand{\arraystretch}{1.8} %mais linhas
 
\begin{tabularx}{\textwidth}{
>{\centering\arraybackslash}p{3cm} % tamanho Variável
>{\raggedright\arraybackslash}X     % Tamanho Descrição
>{\centering\arraybackslash}p{4cm}  % Tamanho Tipo
>{\centering\arraybackslash}p{4cm}  % Tamanho Valores
}
\toprule
\textbf{Variavel} &
\textbf{Descricao} &
\textbf{Tipo} &
\textbf{Valores / unidade} \\
\midrule
 
 
NU\_INSCRICAO&Número da inscrição do aluno&Qualitativa nominal& Número de 12 dígitos \\[0.3cm]
NU\_ANO&Ano de realização da prova&Quantitativa discreta& Anos de 2021 à 2025 \\[0.3cm]
TP\_COR\_RACA& Identificação de raca cor do aluno& Qualitativa nominal & Categorias de Raça/Cor & \\[0.3cm]
NU\_NOTA\_CN& Nota de ciências naturais & Quantitativa contínua & 0 a 1000 pontos \\[0.3cm]
NU\_NOTA\_CH& Nota de ciências humanas& Quantitativa contínua & 0 a 1000 pontos \\[0.3cm]
NU\_NOTA\_LC& Nota de linguagens e codigos & Quantitativa contínua & 0 a 1000 pontos \\[0.3cm]
NU\_NOTA\_MT& Nota de matemática &Quantitativa contínua & 0 a 1000 pontos \\[0.3cm]
NU\_NOTA\_RED& Nota da redação & Quantitativa contínua & 0 a 1000 pontos \\[0.3cm]
Q006& O aluno possui renda? & Quantitativa nominal & A = Não; B = Sim  \\[0.3cm]
Q007& Qual a renda mensal da família do aluno? & Qualitativa ordinal & Valores A < Q \\[0.3cm]
Q020& Na casa do aluno, existe acesso à internet por rede wi-fi? & Qualitativa nominal & A = Não; B = Sim \\[0.3cm]
Q021& Na casa do aluno, existe computador/notebook? & Qualitativa nominal & Categorias de A = Não até E = Sim, quatro ou mais \\[0.3cm]

\bottomrule
\end{tabularx}
 
\end{table}
 
\end{document}
